\clearpage
\VerdierPhysicalPageAnchor{0268}%
\begingroup
\small
\setlength{\parindent}{0pt}%
\setlength{\parskip}{0pt}%
\setlength{\emergencystretch}{3em}%
\newcommand{\VerdierContentsRow}[3]{%
  \par\noindent
  \hangindent=#1\relax\hangafter=1
  \hspace*{#1}#2\nobreak
  \leaders\hbox{\kern.32em.\kern.32em}\hfill
  \nobreak\hbox{\kern.65em #3}\par}
\newcommand{\VerdierContentsChapter}[1]{%
  \par\addvspace{.55\baselineskip}\noindent\textbf{#1}\par\smallskip}
\newcommand{\VerdierContentsSection}[2]{%
  \par\addvspace{.45\baselineskip}\noindent
  \hfill\textbf{#1.}\enspace #2\hfill\mbox{}\par\smallskip}

\vspace*{4.25\baselineskip}
\VerdierSemanticAnchor{toc-title}%
\begin{center}
  {\Large\bfseries Table des matières.}
\end{center}
\vspace{2.5\baselineskip}

\VerdierSemanticAnchor{toc-introduction}%
\VerdierContentsRow{0pt}{\textbf{Introduction}}{1}
\VerdierSemanticAnchor{toc-acknowledgments}%
\VerdierContentsRow{0pt}{\textbf{Remerciements}}{15}

\VerdierSemanticAnchor{toc-ch1}%
\VerdierContentsChapter{Chapitre I : Les catégories de complexes des catégories additives.}
\VerdierSemanticAnchor{toc-ch1-sec1}%
\VerdierContentsSection{1}{Catégories graduées.}
\VerdierSemanticAnchor{toc-ch1-sec1-1}%
\VerdierContentsRow{.5em}{\textbf{1.1.}\enspace Catégories graduées de type $G$}{17}
\VerdierSemanticAnchor{toc-ch1-sec1-2}%
\VerdierContentsRow{.5em}{\textbf{1.2.}\enspace Opération à isomorphisme près d'un groupe sur une catégorie}{19}
\VerdierSemanticAnchor{toc-ch1-sec1-3}%
\VerdierContentsRow{.5em}{\textbf{1.3.}\enspace Équivalence de deux langages}{22}
\VerdierSemanticAnchor{toc-ch1-sec1-4}%
\VerdierContentsRow{.5em}{\textbf{1.4.}\enspace Changement de groupes}{27}
\VerdierSemanticAnchor{toc-ch1-sec1-5}%
\VerdierContentsRow{.5em}{\textbf{1.5.}\enspace $\mathbb Z^n$-catégories}{34}
\VerdierSemanticAnchor{toc-ch1-sec1-6}%
\VerdierContentsRow{.5em}{\textbf{1.6.}\enspace Cocycle de Koszul. Règle des signes}{43}

\VerdierSemanticAnchor{toc-ch1-sec2}%
\VerdierContentsSection{2}{Complexes d'une catégorie additive.}
\VerdierSemanticAnchor{toc-ch1-sec2-1}%
\VerdierContentsRow{.5em}{\textbf{2.1.}\enspace Complexes $n$-uples}{49}
\VerdierSemanticAnchor{toc-ch1-sec2-2}%
\VerdierContentsRow{.5em}{\textbf{2.2.}\enspace Complexe simple associé}{50}
\VerdierSemanticAnchor{toc-ch1-sec2-3}%
\VerdierContentsRow{.5em}{\textbf{2.3.}\enspace Structure graduée sur les catégories de complexes}{54}
\VerdierSemanticAnchor{toc-ch1-sec2-4}%
\VerdierContentsRow{.5em}{\textbf{2.4.}\enspace Extension des foncteurs aux complexes}{56}
\VerdierSemanticAnchor{toc-ch1-sec2-5}%
\VerdierContentsRow{.5em}{\textbf{2.5.}\enspace Homotopies}{60}

\VerdierSemanticAnchor{toc-ch1-sec3}%
\VerdierContentsSection{3}{Triangles distingués.}
\VerdierSemanticAnchor{toc-ch1-sec3-1}%
\VerdierContentsRow{.5em}{\textbf{3.1.}\enspace Le cône d'un morphisme}{65}
\VerdierSemanticAnchor{toc-ch1-sec3-2}%
\VerdierContentsRow{.5em}{\textbf{3.2.}\enspace Propriétés du foncteur cône}{69}
\VerdierSemanticAnchor{toc-ch1-sec3-3}%
\VerdierContentsRow{.5em}{\textbf{3.3.}\enspace Propriétés des triangles distingués}{78}
\VerdierSemanticAnchor{toc-ch1-sec3-4}%
\VerdierContentsRow{.5em}{\textbf{3.4.}\enspace Propriétés fonctorielles des triangles distingués}{84}
\VerdierSemanticAnchor{toc-ch1-appendix}%
\VerdierContentsRow{0pt}{\textbf{Appendice : Commentaires sur le choix des signes}}{91}
\endgroup
